% =====================================================================
% PART 1 (frames 1-6): title, the problem, what was known, main results,
% all k at a glance, roadmap.
% Figures: figs/p1-cutexample.tex, figs/p1-landscape.tex, figs/p1-bounds.tex
% =====================================================================

% local helpers of part 1 (names are prefixed with pI to avoid clashes with other parts)
\providecommand{\pInum}[1]{\tikz[baseline=(n.base)]\node[circle,fill=acc,text=white,font=\bfseries\large,
  inner sep=0pt,minimum size=8.5mm] (n) {#1};}
\providecommand{\pIto}{\ensuremath{\;\longrightarrow\;}}

% ---- Frame 1: title frame
\begin{frame}[plain]
\vspace*{12mm}
{\color{acc}\bfseries\LARGE Improved NP-Hardness of Approximation\\[1.5mm]
 for Max-3-Cut and for Max-$k$-Cut with Small $k$\par}
\vspace{3mm}
{\color{acc!35}\rule{\linewidth}{0.8pt}}\par
\vspace{7mm}
{\Large Ziang Ni\par}
\vspace{2mm}
{\large Delft University of Technology\par}
\vspace{11mm}
{\color{soft}\small STOC 2027 submission video\par}
\end{frame}

% ---- Frame 2: The problem: Max-k-Cut
\begin{frame}{The problem: Max-$k$-Cut}
\begin{columns}[T,onlytextwidth]
\begin{column}{0.53\textwidth}
\raggedright
\vspace{6mm}
\textbf{Input:} graph with non-negative edge weights\par
\vspace{4mm}
\textbf{Colour} every vertex with one of $k$ colours\par
\vspace{4mm}
\textbf{Maximise} the total weight of \hl{cut} edges:\\
edges whose endpoints get different colours\par
\vspace{6.5mm}
{\color{acc!35}\rule{0.94\linewidth}{0.5pt}}\par
\vspace{3mm}
\begin{itemize}\setlength{\itemsep}{2.5mm}
\item $k=2$:\ \ Max-Cut
\item random colouring: each edge is cut with\\ probability $1-\frac1k$\quad ($\frac23$ for $k=3$)
\end{itemize}
\end{column}
\begin{column}{0.45\textwidth}
\centering
\vspace{2mm}
% Frame 2: a small graph with a 3-colouring in which exactly two edges are not cut.
% 8 vertices, 14 unit edges: two K4's that share the vertex c, plus a vertex f joined to a and e.
% Colours: c0 = white, c1 = grey, c2 = accent.  The edges a-b and d-y are monochromatic; every
% 3-colouring of this graph leaves at least two edges uncut (two edge-disjoint K4's), so the
% colouring shown is optimal (checked by brute force: max 3-cut = 12 of 14).
\begingroup
\begin{tikzpicture}[x=0.97cm,y=0.97cm,
  v/.style={circle,draw=black,line width=0.7pt,minimum size=5.8mm,inner sep=0pt},
  c0/.style={v,fill=white},
  c1/.style={v,fill=black!35},
  c2/.style={v,fill=acc},
  cutedge/.style={line width=0.7pt,black},
  uncut/.style={line width=2.2pt,warn,dash pattern=on 3.4pt off 2pt}]
  % vertices
  \node[c0] (a) at (0.20,3.70) {};
  \node[c0] (b) at (0.55,0.50) {};
  \node[c1] (x) at (1.45,2.20) {};
  \node[c2] (c) at (3.10,1.90) {};
  \node[c1] (d) at (5.50,0.30) {};
  \node[c0] (e) at (6.10,3.30) {};
  \node[c1] (y) at (4.90,1.90) {};
  \node[c2] (f) at (3.20,4.30) {};
  % cut edges
  \draw[cutedge] (a)--(x);
  \draw[cutedge] (b)--(x);
  \draw[cutedge] (a)--(c);
  \draw[cutedge] (b)--(c);
  \draw[cutedge] (x)--(c);
  \draw[cutedge] (c)--(d);
  \draw[cutedge] (c)--(e);
  \draw[cutedge] (d)--(e);
  \draw[cutedge] (c)--(y);
  \draw[cutedge] (e)--(y);
  \draw[cutedge] (a)--(f);
  \draw[cutedge] (f)--(e);
  % edges that are not cut (both endpoints have the same colour)
  \draw[uncut] (a)--(b);
  \draw[uncut] (d)--(y);
  % legend
  \draw[cutedge] (0.75,-0.75)--(1.75,-0.75);
  \node[anchor=west,font=\small] at (1.85,-0.75) {cut};
  \draw[uncut] (3.25,-0.75)--(4.25,-0.75);
  \node[anchor=west,font=\small,text=warn] at (4.35,-0.75) {not cut};
\end{tikzpicture}%
\endgroup

\end{column}
\end{columns}
\end{frame}

% ---- Frame 3: Max-3-Cut: what was known
\begin{frame}{Max-3-Cut: what was known}
\vspace{2mm}
Which approximation ratios can be achieved in polynomial time?
\vfill
\centering
% Frame 3: approximation ratios for Max-3-Cut on a linear axis (broken between 0.68 and 0.82).
% Positions (unit 1mm): left piece  x = (a-0.65)*600 for 0.65 <= a <= 0.68,
%                       right piece x = 23 + (a-0.82)*600 for 0.82 <= a <= 1   (same scale: 6mm per 0.01).
% Values: 2/3 = 0.666667, 0.836 (SDP rounding), 27/29 = 0.931034, 16/17 = 0.941176, 55/57 = 0.964912.
% Overlay: everything inside \onslide<2->{...} appears on the second step (the bound of this paper).
\begingroup
\pgfmathsetmacro{\pIxr}{(0.666667-0.65)*600}%      2/3
\pgfmathsetmacro{\pIxs}{23+(0.836-0.82)*600}%      0.836
\pgfmathsetmacro{\pIxn}{23+(0.931034-0.82)*600}%   27/29
\pgfmathsetmacro{\pIxo}{23+(0.941176-0.82)*600}%   16/17
\pgfmathsetmacro{\pIxg}{23+(0.964912-0.82)*600}%   55/57
\pgfmathsetmacro{\pIxe}{23+(1-0.82)*600}%          1
\begin{tikzpicture}[x=1mm,y=1mm,
  declare function={axa(\a)=(\a-0.65)*600; axb(\a)=23+(\a-0.82)*600;},
  lab/.style={font=\small,inner sep=1pt,align=left},
  reg/.style={font=\normalsize,inner sep=0pt,align=center}]
  \useasboundingbox (-3,-29) rectangle (139,31.5);
  % ---- the three regions (bars, 0 <= y <= 14) ----
  \fill[black!7]  (0,0) rectangle (\pIxs,14);
  \fill[black!25] (\pIxo,0) rectangle (\pIxe,14);
  \onslide<2->{\fill[acc] (\pIxn,0) rectangle (\pIxo,14);}
  \draw[black!60,line width=0.5pt] (0,14)--(\pIxe,14);
  \draw[black!60,line width=0.5pt] (0,0)--(0,14) (\pIxe,0)--(\pIxe,14);
  \draw[black!60,line width=0.5pt] (\pIxs,0)--(\pIxs,14);
  \node[reg] at ({\pIxs/2},7) {polynomial\\[-1pt]time};
  \node[reg] at ({(\pIxs+\pIxn)/2},7)
     {\textbf{open}\\[-0.5pt]{\small\color{soft}(assuming only $\mathrm{P}\neq\mathrm{NP}$)}};
  \node[reg] at ({(\pIxo+\pIxe)/2},7) {\textbf{NP-hard}};
  % ---- axis with a break (no tick marks at 0.68 and 0.82: they would touch the break marks) ----
  \draw[line width=0.9pt] (0,0)--(19.3,0) (21.7,0)--(\pIxe,0);
  \draw[line width=0.7pt] (18.5,-1.6)--(20.1,1.6) (20.9,-1.6)--(22.5,1.6);
  \foreach \a in {0.65,0.66,0.67} \draw[black!70,line width=0.4pt] ({axa(\a)},0)--({axa(\a)},-1.1);
  \foreach \a in {0.83,0.84,...,1.001} \draw[black!70,line width=0.4pt] ({axb(\a)},0)--({axb(\a)},-1.1);
  \foreach \a/\t in {0.9/0.9,1/1}
    {\draw[black!70,line width=0.5pt] ({axb(\a)},0)--({axb(\a)},-2);
     \node[font=\small,text=soft,anchor=north,inner sep=1pt] at ({axb(\a)},-2.8) {$\t$};}
  % ---- algorithms (below the axis) ----
  \draw[line width=0.9pt] (\pIxr,0)--(\pIxr,-2.4);
  \node[lab,anchor=north,align=center] at (\pIxr,-2.8) {$\mathbf{2/3}$};
  \node[lab,anchor=north,align=center] at (\pIxr,-9.4) {random\\[-1pt]colouring};
  \draw[line width=0.9pt] (\pIxs,0)--(\pIxs,-2.4);
  \node[lab,anchor=north] at (\pIxs,-2.8) {$\mathbf{0.836}$};
  \node[lab,anchor=north west] at ({\pIxs-5.2},-9.4)
     {best algorithm: SDP rounding \src{FJ97, GW04, dKPW04}\\[0.5pt]
      optimal under the Unique Games Conjecture};
  % ---- [NRT25] for k = 3 (grey) ----
  \fill[black!55] (\pIxg,0) circle[radius=0.9];
  \draw[black!55,line width=0.6pt] (\pIxg,0)--(\pIxg,-13.6);
  \node[lab,anchor=north west,text=soft] at ({\pIxg-0.2},-14.2)
     {$55/57\approx 0.965$\\[0.5pt]gadgets of [NRT25]\\[-1pt](AlphaEvolve)};
  % ---- hardness (above the bars) ----
  \draw[line width=0.9pt] (\pIxo,0)--(\pIxo,20.5);
  \node[lab,anchor=south west] at ({\pIxo+1.2},15.6)
     {$\mathbf{16/17\approx 0.941}$\\[0.5pt]NP-hard \src{AOTW14}, 2012};
  % ---- step 2: this paper ----
  \onslide<2->{
    \draw[acc,line width=1.1pt] (\pIxn,0)--(\pIxn,20.5);
    \node[lab,anchor=south east,align=right,text=acc] at ({\pIxn-1.2},15.6)
       {$\mathbf{27/29\approx 0.931}$\\[0.5pt]\textbf{NP-hard: this paper}};
    \draw[acc,line width=1.1pt,-{Stealth[length=2mm,width=2mm]}] ({\pIxo-0.5},17.6)--({\pIxn+0.6},17.6);
  }
\end{tikzpicture}%
\endgroup

\vfill
\end{frame}

% ---- Frame 4: Main results
\begin{frame}{Main results}
\vspace{-0.5mm}
\begin{block}{Theorem 1.1\ \ (Max-3-Cut)}
\begin{minipage}[c]{0.545\linewidth}
For every $\eps>0$ it is NP-hard to distinguish\\[0.4mm]
\hspace*{3mm}\textbf{some} $3$-colouring cuts $\ge\frac{29}{30}-\eps$\\[0.4mm]
\hspace*{3mm}\textbf{every} $3$-colouring cuts $\le\frac{9}{10}+\eps$\\[0.2mm]
{\footnotesize\color{soft}\hspace*{3mm}(as fractions of the total edge weight)}
\end{minipage}\hfill
\begin{minipage}[c]{0.44\linewidth}
\centering
$\Longrightarrow$\ \ NP-hard to approximate within\\[1.6mm]
{\Large\hl{$\mathbf{27/29}+\eps\approx0.931$}}\\[1.6mm]
{\small\color{soft}previous: $16/17\approx0.941$} \src{AOTW14}
\end{minipage}
\end{block}
\vspace{-0.5mm}
\begin{columns}[T,onlytextwidth]
\begin{column}{0.475\textwidth}
\begin{block}{Theorem 1.2\ \ (every $k\ge3$)}
NP-hard to approximate within\\[0.4mm]
$1-2/\Gamma_{k-3}+\eps$,\ \ \ $\Gamma_L=29+18L+3\tbinom L2$\\[1.5mm]
\begin{tabular}{@{}l@{\hspace{6mm}}l@{}}
$k=4$ & $k=5$\\[-0.3mm]
\hl{$45/47\approx0.957$} & \hl{$33/34\approx0.971$}\\[-0.6mm]
{\footnotesize\color{soft}was $0.987$} \src{NRT25} & {\footnotesize\color{soft}was $120/121$} \src{AOTW14}
\end{tabular}
\end{block}
\end{column}
\begin{column}{0.482\textwidth}
\begin{block}{Theorem 1.3\ \ ($k$-colourable graphs)}
NP-hard to distinguish: $k$-colourable, or\\[0.4mm]
every $k$-colouring cuts at most\vphantom{$\tbinom L2$}\\[1.5mm]
\begin{tabular}{@{}l@{\hspace{3mm}}l@{}}
$k=4$ & $k=5$\\[-0.3mm]
\hl{$82/85+\eps\approx0.965$} & \hl{$119/122+\eps\approx0.975$}\\[-0.6mm]
{\footnotesize\color{soft}stated before: $85/86$} & {\footnotesize\color{soft}stated before: $120/121$}
\end{tabular}
\end{block}
\end{column}
\end{columns}
\end{frame}

% ---- Frame 5: All k at a glance, and one limitation
\begin{frame}{All $k$ at a glance, and one limitation}
\begin{columns}[c,onlytextwidth]
\begin{column}{0.62\textwidth}
% Frame 5: the paper's Figure 1 (paper/figures/bounds.tex) adapted for a slide.
% ALL plot coordinates, the axis limits, the y ticks and the shaded range are copied unchanged from
% paper/figures/bounds.tex (compared by script).  Presentation changes only: size (83mm x 68mm), larger
% fonts, slightly thicker lines, hard-coded legend texts (no \ref / \cite in the slides), x tick labels
% only at 3,5,10,15,20,24 (tick marks at every k as before), x label "number of colours k", y label note
% "(higher is stronger)", shorter note in the shaded range, colour `acc` (= blue!60!black of the paper).
%   Thm 1.2:  2/Gamma_L,                    Gamma_L = 29 + 18L + 3 binom(L,2),  L = k-3
%   Thm 1.3:  3/(51 + 34L + 3 binom(L,2)),  k >= 4
%   [AOTW14]: 1/(17(k+c_k)+c_k), c_k = k mod 3, for k >= 4;  1/17 for k = 3
%   [NRT25]:  1 - 0.98708 = 0.01292 at k = 4
\begingroup
%
\begin{tikzpicture}
\begin{axis}[
  width=83mm, height=68mm,
  scale only axis=false,
  xmin=2.5, xmax=24.5,
  ymode=log, ymin=0.0016, ymax=0.09,
  xtick={3,5,10,15,20,24},
  minor xtick={3,...,24},
  ytick={0.002,0.005,0.01,0.02,0.05},
  yticklabels={$0.002$,$0.005$,$0.01$,$0.02$,$0.05$},
  minor ytick={0.003,0.004,0.006,0.007,0.008,0.009,0.03,0.04,0.06,0.07,0.08,0.09},
  ymajorgrids, yminorgrids=false, xmajorgrids=false, xminorgrids=false,
  major grid style={black!14,line width=0.3pt},
  tick label style={font=\footnotesize},
  tick style={black!70},
  xlabel={number of colours $k$}, xlabel style={font=\small,yshift=3pt},
  ylabel={$1-{}$hardness factor\\[-2pt]{\footnotesize(higher is stronger)}},
  ylabel style={font=\small,align=center,yshift=-3pt},
  axis on top,
  legend cell align=left,
  legend pos=north east,
  legend style={font=\scriptsize,draw=black!40,fill=white,inner sep=2pt,
                row sep=-1pt,line width=0.4pt},
  every axis plot/.append style={line join=round},
]
% range of k in which this paper improves the previous bounds
\fill[black!6] (axis cs:2.5,0.0016) rectangle (axis cs:20.5,0.09);
\node[font=\scriptsize,text=black!60,anchor=south west,inner sep=0pt]
  at (axis cs:3.3,0.00207) {shaded: improved by this paper};
% (1) Theorem 1.2
\addplot[acc, line width=1pt, mark=*, mark size=1.5pt,
         mark options={solid,fill=acc}]
  coordinates {(3,0.0689655) (4,0.0425532) (5,0.0294118) (6,0.0217391) (7,0.0168067)
   (8,0.0134228) (9,0.010989) (10,0.00917431) (11,0.0077821) (12,0.00668896)
   (13,0.00581395) (14,0.00510204) (15,0.00451467) (16,0.00402414) (17,0.00361011)
   (18,0.00325733) (19,0.00295421) (20,0.00269179) (21,0.00246305) (22,0.00226244)
   (23,0.00208551) (24,0.00192864)};
\addlegendentry{this paper (Thm 1.2)}
% (2) Theorem 1.3
\addplot[black, line width=0.7pt, dashed, mark=square*, mark size=1.3pt,
         mark options={solid,fill=white}]
  coordinates {(4,0.0352941) (5,0.0245902) (6,0.0185185) (7,0.0146341) (8,0.0119522)
   (9,0.01) (10,0.00852273) (11,0.00737101) (12,0.00645161) (13,0.00570342)
   (14,0.00508475) (15,0.00456621) (16,0.00412655) (17,0.00375) (18,0.00342466)
   (19,0.00314136) (20,0.00289296) (21,0.0026738) (22,0.00247934) (23,0.00230592)
   (24,0.00215054)};
\addlegendentry{this paper, $k$-colourable (Thm 1.3)}
% (3) previous bound of [AOTW14]
\addplot[black!65, line width=0.7pt, densely dotted, mark=triangle*, mark size=1.7pt,
         mark options={solid,fill=black!65}]
  coordinates {(3,0.0588235) (4,0.0116279) (5,0.00826446) (6,0.00980392) (7,0.00729927)
   (8,0.00581395) (9,0.00653595) (10,0.00531915) (11,0.0044843) (12,0.00490196)
   (13,0.0041841) (14,0.00364964) (15,0.00392157) (16,0.00344828) (17,0.00307692)
   (18,0.00326797) (19,0.00293255) (20,0.00265957) (21,0.00280112) (22,0.00255102)
   (23,0.00234192) (24,0.00245098)};
\addlegendentry{previous [AOTW14]}
% (4) [NRT25] at k = 4
\addplot[only marks, black, mark=diamond*, mark size=2.1pt, mark options={fill=white}]
  coordinates {(4,0.01292)};
\addlegendentry{[NRT25]}
\end{axis}
\end{tikzpicture}%
\endgroup

\end{column}
\begin{column}{0.345\textwidth}
\begin{itemize}\setlength{\itemsep}{3.5mm}
\item \hl{improved} for every $k$ from $3$ to $20$
\item $k\ge21$: previous bound\\ \src{AOTW14} remains best
\end{itemize}
\vspace{6mm}
\begin{alertblock}{Limitation}
No perfect completeness:\\ on $3$-colourable graphs,\\ $16/17$ \src{AOTW14} is still\\ the best bound.
\end{alertblock}
\end{column}
\end{columns}
\end{frame}

% ---- Frame 6: Three ideas
\begin{frame}{Three ideas}
\vspace{2mm}
\renewcommand{\arraystretch}{1.0}
\begin{tabular}{@{}m{12mm}@{\hspace{3mm}}m{125mm}@{}}
\pInum{1} & {\Large\textbf{The negation trick}\pIto\hl{$\mathbf{27/29}$ for Max-3-Cut}}\\[1mm]
          & {\color{soft}pair every assignment $X$ with $-X$; proof in a few lines}
            \hfill\paperref{Section 3}\\[10mm]
\pInum{2} & {\Large\textbf{A soundness criterion} for gadget reductions\pIto\hl{all $k$}}\\[1mm]
          & {\color{soft}from H\aa stad's equations; used for more than three colours}
            \hfill\paperref{Sections 4--5}\\[10mm]
\pInum{3} & {\Large\textbf{A barrier}\pIto\hl{$\mathbf{27/29}$ is optimal for this method}}\\[1mm]
          & {\color{soft}under the criterion, gadgets of any size cannot do better}
            \hfill\paperref{Section 7}\\[2.5mm]
\multicolumn{2}{@{}l@{}}{\small\color{soft}At the end: how these results were obtained --- with extensive use of an AI system}
\end{tabular}
\end{frame}
